\documentclass{article}
\usepackage{amsmath,amssymb}
\title{A grouped revision example}
\author{}
\date{}
\begin{document}
\maketitle

\section{Assumptions}
Let \(f:\mathbb R\to\mathbb R\) be differentiable at \(x\).
We consider positive increments \(h\).
The derivative describes the local first-order change in the function.

\section{Local approximation}
The approximation is written as an exact equality:
\[
  f(x+h)=f(x)+f'(x)h.
\]
The remainder is omitted from this draft.

\section{A worked example}
Take \(f(t)=t^2\) and keep \(x\) fixed.
Then \(f'(x)=x\).

\section{Temporary checklist}
This section is only a drafting reminder.
Check the assumptions before sharing the note.
Add a worked calculation if there is room.
Remove this checklist before publication.

\section{Reading the result}
A useful explanation names the assumptions, keeps each calculation near the claim it supports, and states which terms disappear in the limit so that readers can check the reasoning without jumping between distant paragraphs.
The next version will contain a worked calculation.

\end{document}
